% TOP_PROBLEM: {"id":30007687,"problem_number":"LOCAL-30007687","title":"Reputation versus public regulation","category_id":21,"category":{"id":21,"name":"economics","display_name":"Economics","description":"Open economic theory statements, published research agendas, and proposed scoped specifications; question provenance and historical priority are qualified per record.","slug":"economics","order_index":21,"created_at":"2026-10-08T00:00:00Z"},"status":"research_question","external_url":null,"legacy_ids":[],"published":false,"statement":"Compare platform reputation, public inspection and their combination for repair-service welfare and independently measured safety.","economics":{"original_id":176,"field":"Digital, platform and data economics","kind":"empirical","formalization_basis":"adapted_research_question","source_status":"research_agenda","priority_status":"earliest_found","priority_note":"Earliest relevant explicit question or agenda passage located in this audit. No claim of first-ever formulation; earlier versions and later solutions were not exhaustively traced.","earliest_verified_reference":{"authors":["Chiara Farronato"],"title":"Digital Platforms as Information Aggregators: Implications for their Design and Regulation","year":2025,"url":"https://conference.nber.org/confer/2025/DTs25/farronato.pdf","doi":null,"locator":"Section 2.1, printed pp. 7–8, platform monitoring and public regulation","evidence_summary":"Calls for linked quality data to evaluate substitution or complementarity with regulation."},"current_reference":{"authors":["Chiara Farronato"],"title":"Digital Platforms as Information Aggregators: Implications for their Design and Regulation","year":2025,"url":"https://conference.nber.org/confer/2025/DTs25/farronato.pdf","doi":null,"locator":"Section 2.1, printed pp. 7–8, platform monitoring and public regulation","evidence_summary":"Calls for linked quality data to evaluate substitution or complementarity with regulation."},"formalization_note":"The source verifies a broader question or research agenda; the population, estimand, equations and restrictions here are our scoped adaptation. This is not a quoted canonical conjecture, and the audit does not prove contemporary unsolved status for every special case."},"statusQualification":"Published question or research agenda verified at the cited locator. The mathematical specification is adapted for this catalogue and includes additional assumptions; source publication does not establish first-ever priority or certify this exact model remains unsolved. The source verifies a broader question or research agenda; the population, estimand, equations and restrictions here are our scoped adaptation. This is not a quoted canonical conjecture, and the audit does not prove contemporary unsolved status for every special case.","statusReviewedAt":"2026-10-08"}
% FIRST_POSTED: 2026-10-08
% ENTRY_KIND: statement_only
% !TeX program = lualatex
\documentclass[11pt]{article}
\usepackage{fontspec}
\usepackage[a4paper,margin=25mm]{geometry}
\usepackage{amsmath,amssymb,mathtools}
\usepackage{xurl}
\usepackage[hidelinks]{hyperref}
\newcommand{\sref}[2]{\hyperlink{source:#1:#2}{[S#2]}}
\setlength{\emergencystretch}{3em}
\begin{document}
% problemId:30007687
\section*{Reputation versus public regulation}\label{30007687}
\subsection{Definitions and mathematical statement}
Consider home-repair providers serving households in a fixed set of United States cities over three years. Let \(r\in\{0,1,2\}\) denote regulation by licensing and inspection alone, platform reputation alone, or a specified combination. Under each regime providers may enter, exit and change service effort. Let \(H_i(r)\) be household \(i\)'s independently audited safety loss, including latent defects discovered within three years, and \(V_i(r)\) its money-valued service benefit minus payment and search cost. Let \(K(r)\) be real public and platform monitoring cost per eligible household. Define \(W(r)=E[V_i(r)-H_i(r)]-K(r)\), where expectation covers households needing a repair, including those unable to hire under the regime. The empirical task is to compare \(W(r)\) and the distribution of \(H_i(r)\), not merely ratings among completed transactions. Use regulatory changes or market-level randomized monitoring with explicit counterfactual assumptions, and link platform histories to independent safety inspections. Quantify whether reputation observes the dimensions of quality protected by public regulation, and whether cheaper screening changes provider composition. Account for selective review submission and households switching between online and offline suppliers. A substitution claim requires both welfare gains and compliance with prespecified safety thresholds, with uncertainty about unobserved defects reported separately.

\par\textbf{Formulation and scope.} The source verifies a broader question or research agenda; the population, estimand, equations and restrictions here are our scoped adaptation. This is not a quoted canonical conjecture, and the audit does not prove contemporary unsolved status for every special case.

\par\textbf{Open-question provenance.} An explicit question or research agenda was verified at the source locator. This does not by itself certify that every later version remains unresolved \sref{30007687}{1}.

\par\textbf{Historical priority.} Earliest relevant explicit question or agenda passage located in this audit. No claim of first-ever formulation; earlier versions and later solutions were not exhaustively traced.

\subsection{Short English statement}
Compare platform reputation, public inspection and their combination for repair-service welfare and independently measured safety.

\subsection{Sources}
\begin{itemize}\raggedright
\item[\textnormal{[S1]}]\hypertarget{source:30007687:1}{} Chiara Farronato. 2025. Digital Platforms as Information Aggregators: Implications for their Design and Regulation. Section 2.1, printed pp. 7–8, platform monitoring and public regulation.
\url{https://conference.nber.org/confer/2025/DTs25/farronato.pdf}.
{\small\textit{Earliest verified question/agenda reference; historical priority qualified below. Calls for linked quality data to evaluate substitution or complementarity with regulation.}}
\end{itemize}
\end{document}
